\documentclass[answers,12pt,addpoints]{exam} 
\usepackage{import}

\import{C:/Users/prana/OneDrive/Desktop/MathNotes}{style.tex}

% Header
\newcommand{\name}{Pranav Tikkawar}
\newcommand{\course}{16:640:654}
\newcommand{\assignment}{Stochastic Processes}
\author{\name}
\title{\course \ - \assignment}

\begin{document}
\maketitle

\newpage
\tableofcontents
\newpage
\section{1/21}
\begin{definition}[Stochastic Process]
    A collection of random variables $\{X_t\}_{t \in T}$ defined on a common probability space $(\Omega, \mathcal{F}, P)$, where $T$ is an index set (often representing time). Each random variable $X_t$ maps outcomes in $\Omega$ to a state space $S$.
\end{definition}
\begin{definition}[Random Walk]
    Let $\{X_i\}_{i=1}^{\infty}$ be a sequence of i.i.d. random variables. DEfine the random walk $\{S_n\}_{n=0}^{\infty}$ by
    \[
        S_n = \sum_{i=1}^{n} X_i, \quad S_0 = 0.
    \]
    Then $\{S_n\}$ is called a random walk.

    Can do random walks on $\mathbb{Z}, \mathbb{R}, \mathbb{R}^d$, as well as graphs. In particular, if $G = \mathbb{Z}^d$ random walk is a d-dim random walk.

    $P(X_i = 1) = p$, $P(X_i = -1) = q = 1-p$ is called a simple random walk on $\mathbb{Z}$ and if $p = q = \frac{1}{2}$, it is called a symmetric random walk.
\end{definition}
\begin{example}[Simple Random Walk on $\mathbb{Z}^d$]
    A simple random walk on $\mathbb{Z}^d$ starts at $vec(0)$ and each step given the current state $vec{x} \in \mathbb{Z}^d$, moves to one of the $d$ neighboring states $vec{x} + vec{e_i}$ or $vec{x} - vec{e_i}$ with probability $\frac{1}{2}$. 

    If we define $\tau(0) = \inf\{t \geq 1: S_n = 0\}$, we say that the random walk is recurrent if $P(\tau(0) < \infty) = 1$ and transient if $P(\tau(0) < \infty) < 1$.
\end{example}
\begin{theorem}[Polya's Recurrence Theorem]
    A simple random walk on $\mathbb{Z}^d$ is recurrent for $d = 1, 2$ and transient for $d \geq 3$.
\end{theorem}
\begin{note}[Stiens Phenomenon]
    Given data $x_1, x_2, \ldots, x_n \in \mathbb{R}^d$, iid $\sim N(\theta, I)$ when $d > 3$, the Mean estimator $\hat{\theta} = \bar{x}$ is inadmissible under squared error loss. IE there is a strictly better estimator. This is called Steins Phenomenon.
\end{note}
\begin{definition}[Martingale]
    Example (Betting Game): Suppose each round you have win prob $p$ so $\{ \psi_1, \psi_2, \ldots \}$ are iid bernoulli($p$) random variables. Assuming betting $1$  each round then $X_i = 2\psi_i - 1$ is the gamblers gain in round $i$. Then the net gain after $n$ rounds is a random walk $S_n = \sum_{i=1}^{n} X_i$. When $p = \frac{1}{2}$, the game is fair and $E[X_i] = 0$. the game is fair and $S_n$ is 1-d symmetric random walk. This is the simplest example of a martingale.
    
    If $p < \frac{1}{2}$, the game is unfavorable and called a submartingale. If $p > \frac{1}{2}$, the game is favorable and called a supermartingale.

    If stop betting when you reach a certain amount of money, winner or loser, then we can ask what the expected stopping time, what is its distribution, etc.

    Also we can ask about what would be our betting strategy to maximize our expected gain.
\end{definition}
\begin{questions}
    \question Suppose we toss a fair coin, what is the expected time until we see $HTT$?
    \begin{solution}
        We can model this as a Markov Chain with states:
        \begin{itemize}
            \item $S_0$: No matches
            \item $S_1$: Matched H
            \item $S_2$: Matched HT
            \item $S_3$: Matched HTT (absorbing state)
        \end{itemize}
        Then the probability transition matrix is
        \begin{align*}
            \begin{bmatrix}
                0.5 & 0.5 & 0 & 0 \\
                0 & 0.5 & 0.5 & 0 \\
                0.5 & 0 & 0 & 0.5 \\
                0 & 0 & 0 & 1
            \end{bmatrix}
        \end{align*}
        Let $E_i$ be the expected time to reach $S_3$ from state $S_i$. Then we have the following equations:
        \begin{align*}
            E_0 &= 1 + 0.5E_0 + 0.5E_1 \\
            E_1 &= 1 + 0.5E_2 + 0.5E_0 \\
            E_2 &= 1 + 0.5E_3 + 0.5E_0 \\
            E_3 &= 0
        \end{align*}
        Solving these equations, we find:
        \begin{align*}
            E_0 &= 8 \\
            E_1 &= 6 \\
            E_2 &= 4
        \end{align*}
        Therefore, the expected time until we see the sequence $HTT$ is $E_0 = 8$.
    \end{solution}
    \question A monkey is randomly typing on a typewriter with 26 letters. What is the expected time until the monkey types "ABRACADABRA"?
\end{questions}
\begin{definition}[Markov Chain]
    Example (Random walk on finite graph): Let $G = (V, E)$ be a finite, undirected graph. If we define $deg(v)$ as the number of neighbors of vertex $v \in V$. A random walk on $G$ can be described with the following transition matrix:
    \begin{align*}
        \mathbb{P}(u,v) = \begin{cases}
            0 & \text{if } (u,v) \notin E \\
            \frac{1}{deg(u)} & \text{if } (u,v) \in E
        \end{cases}
    \end{align*}
    Then define $\pi(u) = \frac{deg(u)}{2|E|}$ where $|E|$ is the number of edges in the graph. Then $\pi$ is a stationary distribution for the random walk on $G$.
    ie if $X_0 \sim \pi$, then $X_n \sim \pi$ for all $n \geq 0$.
    \begin{proof}
        We first check $\pi(u)P(u,v) = \pi(v)P(v,u)$ for all $u, v \in V$.This is called the detailed balance condition.
        \begin{itemize}
            \item If $(u,v) \notin E$, then both sides are $0$.
            \item If $(u,v) \in E$, then
            \begin{align*}
                \pi(u)P(u,v) &= \frac{deg(u)}{2|E|} \cdot \frac{1}{deg(u)} = \frac{1}{2|E|} \\
                \pi(v)P(v,u) &= \frac{deg(v)}{2|E|} \cdot \frac{1}{deg(v)} = \frac{1}{2|E|}
            \end{align*}
        \end{itemize}
        Thus the detailed balance condition holds. 

        Now we check that $\mathbb{P}(X_1 = u) = \sum_{v \in V} \mathbb{P}(X_1 = u | X_0 = v) \mathbb{P}(X_0 = v) = \sum_{v \in V} P(v,u) \pi(v) = \sum_{v \in V} P(u,v) \pi(u) = \pi(u) \sum_{v \in V} P(u,v) = \pi(u) \cdot 1 = \pi(u)$.

        Thus we need to ensure that $\sum_{v \in V} P(u,v) = 1$, which is true since $P$ is a transition matrix.
    \end{proof}
\end{definition}
\section{1/26}
\begin{definition}[Martingale]
    Is used to model a fair fame. In which you know the history, but the expected future value is equal to the present value.

    $E[Z_n] < \infty$ for all $n$.$E[Z_{n+1} | Z_1, Z_2, \ldots, Z_n] = Z_n$.
\end{definition}
\begin{remark}
    IF we interprt Z as a gamblers fortune after the nth bet the expected fortune is equal to the current fortune.
\end{remark}
\begin{definition}[Formal def. Martingale]
    ${Z_n}$ is a martingale wrt to a filtration ${\mathcal{F}_n}$ if
    \begin{itemize}
        \item $Z_n$ is $\mathcal{F}_n$ measurable
        \item $E[|Z_n|] < \infty$
        \item $E[Z_{n+1} | \mathcal{F}_n] = Z_n$
    \end{itemize}
\end{definition}
\begin{example}[Martingale Example]
    Let $S_n = \sum_n psi_i$ for i.i.d. random variables 
\end{example}
\begin{example}[Quadratic Martingale]
    
\end{example}
\begin{example}[Exponential Martingale]
    
\end{example}
\begin{example}
    Doob's Martingale
\end{example}

\section{1/28}
\begin{remark}
    The central definitin of a martingale is $E[X_{n+1} | \text{past}] = X_n$.

    Recall a probability space: $(\Omega, \mathcal{F}, \mathbb{P})$ 
    
    $\Omega$ is the sample space, $\mathcal{F}$ is a sigma algebra on $\Omega$, and $\mathbb{P}$ is a probability measure on $(\Omega, \mathcal{F})$.
\end{remark}

\begin{definition}[$\sigma$-algebra]
    A collection $\mathcal{F}$ of subsets of $\Omega$ is called a sigma-algebra if:
    \begin{itemize}
        \item $\Omega \in \mathcal{F}$
        \item If $A \in \mathcal{F}$, then $A^c \in \mathcal{F}$
        \item If $A_1, A_2, \ldots \in \mathcal{F}$, then $\bigcup_{i=1}^{\infty} A_i \in \mathcal{F}$
    \end{itemize}
    
\end{definition}
\begin{definition}[Conditional Expectaion]
    Let $(\Omega, \mathcal{F}, \mathbb{P})$ be a probability space, and let $\mathcal{G} \subseteq \mathcal{F}$ be a sub-sigma-algebra. For a random variable $X$ with $E[|X|] < \infty$, the conditional expectation of $X$ given $\mathcal{G}$, denoted by $E[X | \mathcal{G}]$, is a $\mathcal{G}$-measurable random variable satisfying:
    \[
        E[X \cdot 1_A] = E[E[X | \mathcal{G}] \cdot 1_A] \quad \text{for all } A \in \mathcal{G}.
    \]
    
\end{definition}

\section{2/4}

\begin{definition}[Filtration]
    A filtration on a probability space is an incresing sequence of $\sigma$-algebras $\{\mathcal{F}_n\}_{n \geq 0}$ such that $\mathcal{F}_n \subseteq \mathcal{F}_{n+1} \subseteq \mathcal{F}$ for all $n \geq 0$.

    This is typically used to describe the information that is available up to time $n$.
\end{definition}

\begin{definition}[Adapted r.v]
    Given a filtration $\{\mathcal{F}_n\}$, and a stochastic process $\{Z_n\}$, we say that $\{Z_n\}$ is adapted to the filtration if for each $n$, the random variable $Z_n$ is $\mathcal{F}_n$-measurable.
\end{definition}

\begin{definition}[Martingale]
    Given a probability space $(\Omega, \mathcal{F}, \mathbb{P})$ with a filtration $\{\mathcal{F}_n\}$, a stochastic process $\{Z_n\}$ is called a martingale with respect to the filtration if:
    \begin{itemize}
        \item $E|Z_n| < \infty$ for all $n$.
        \item $Z_n$ is adapted to the filtration $\{\mathcal{F}_n\}$.
        \item $E[Z_{n+1} | \mathcal{F}_n] = Z_n$ for all $n$.
    \end{itemize}
    Moreover under 1,2 but not 3 it is called a submartingale if $E[Z_{n+1} | \mathcal{F}_n] \geq Z_n$ and supermartingale if $E[Z_{n+1} | \mathcal{F}_n] \leq Z_n$.
\end{definition}

\begin{definition}[Stopping Time]
    Given a filtration $\{\mathcal{F}_n\}$, a random variable $\tau: \Omega \to \{0, 1, 2, \ldots, \infty\}$ is called a stopping time w.r.t $\{\mathcal{F}_n\}$ if
    $$
        \{\omega : \tau(\omega) = n\} \in \mathcal{F}_n \quad \text{for all } n \geq 0.
    $$
    This means that the event that the stopping time equals $n$ is determined by the information available up to time $n$.
\end{definition}

\begin{lemma}
    Let $T$ be a stopping time with respect to the filtration $\{\mathcal{F}_n\}$. then 
    \begin{itemize}
        \item $\{T \leq n\} \in \mathcal{F}_n$ for all $n$.
        \item $\{T \geq n\} \in \mathcal{F}_{n-1}$ for all $n \geq 1$.
        \item $S = T \wedge m = \min(T, m)$ is a stopping time for any fixed $m$.
    \end{itemize}
\end{lemma}

\begin{lemma}[Walds Lemma]
    Let $\{F_n\}$ be a filtration and $\{Z_n\}$ be a martingale with respect to $\{F_n\}$. And $T$ is a stopping time with respect to $\{F_n\}$ Suppose $T \leq N$ for some constant $N$ almost surely. Then
    $$
        E[Z_T] = E[Z_0].
    $$
    
\end{lemma}


\section{2/9}
\begin{remark}[Convergence Theorem]
    Different types of convergence for random variables:
    \begin{itemize}
        \item Convergence in Probability: $X_n \xrightarrow{P} X$ if for all $\epsilon > 0$, $P(|X_n - X| > \epsilon) \to 0$ as $n \to \infty$.
        \item Convergence in Distribution: $X_n \xrightarrow{d} X$ if the distribution functions of $X_n$ converge to the distribution function of $X$ at all continuity points of the distribution function of $X$.
        \item Almost Sure Convergence: $X_n \xrightarrow{a.s.} X$ if $P(\lim_{n \to \infty} X_n = X) = 1$.
    \end{itemize}
\end{remark}
\begin{definition}[Almost Sure Convergence]
    For a sequence of r.vs $\{X_n\}$ we say that $X_n \to X$ almost surely if $P(\{\omega : \lim_{n \to \infty} X_n(\omega) = X(\omega)\}) = 1$.

    We want to know if $X_n \to X$ a.s when will $E[X_n] \to E[X]$? 

    \begin{example}
        $\Omega = [0,1]$ $P$ is uniform distribution. Define $X_n(\omega) = n!_{[0,1/n]}$ We know that $X_n \to 0$ a.s, but $E[X_n] = 1$ for all $n$ so $E[X_n] \not\to E[X]$.
    \end{example}
\end{definition}
\begin{definition}[Monotone Convergence Theorem]
    If $X_1 \leq X_2 \leq \ldots$ and $X_n \to X$ a.s, then $\lim_{n \to \infty} E[X_n] = E[ \lim_{n \to \infty} X_n]$.

    Ie we can exchange limit and expectation if the sequence is dominated by an integrable random variable.
    
\end{definition}
\begin{definition}[Dominated Convergence Theorem]
    If $X_n \to X$ a.s and $X_n \leq X_{n+1}$ for all $n$, then $E[X_n] \to E[X]$.
        
\end{definition}
\begin{remark}
        Walds equation I says if $\{X_n\}$ is martingale $T$ stopping time, $T \leq N$ for constant $N$ then $E[X_T] = E[X_1]$.

        To apply Walds equation I, we often define, define $T = T \wedge N$ so we have $E[X_{T_n}] = E[X_1]$ then argue that $E[X_{T_n}] \to E[X_T]$ as use DCT to argue that $E[X_{T_n}] \to E[X_T]$ as $n \to \infty$.
\end{remark}
\begin{example}[Simple Random Walk]
    $\{S_n\}$ is a 1D Simple random walk, fix $a,b >0$ and define $T_{a,b} = \inf\{t : S_t = -a \text{ or } S_t = b\}$. 

    Then $\mathbb{P}(T_{a,b} = -a) = \frac{b}{a+b}$ and $\mathbb{P}(T_{a,b} = b) = \frac{a}{a+b}$ and $E[T_{a,b}] = ab$.

    proof of $E[T_{a,b}] = ab$: Recall $S_n^2 - n$ is a martingale 
    \begin{align*}
        E[S_{n+1}^2 - (n+1) | X_1, \ldots, X_n] &= E[(S_n + X_{n+1})^2| X_1, \ldots, X_n] - (n+1) \\
        &= E[S_n^2 + 2S_nX_{n+1} + X_{n+1}^2 | X_1, \ldots, X_n] - (n+1) \\
        &= S_n^2 + 1 + 0 - (n+1) \\
        &= S_n^2 - n
    \end{align*}
    he stopping time for $\{S_n^2 - n\}$ is $S_{T_{a,b}}^2 - T_{a,b}$. Then $E[T_{a,b}] = E[S_{T_{a,b}}^2] = a^2 \cdot \frac{b}{a+b} + b^2 \cdot \frac{a}{a+b} = ab$.

    We define $T_n = T_{a,b} \wedge n$ for each fixed $n$ Apply Wald 1 on $\{S_n^2 - n\}$ at $T_n$ we have 
    \begin{align*}
        E[S_{T_n}^2 - T_n] &= E[S_1^2 - 1] = 0 \\
        E[T_n] &= E[S_{T_n}^2]
    \end{align*}
    For RHS $T_n = T_{a,b} \wedge n$ By MCT we have $\lim_{n \to \infty} E[T_n] = E[T]$ For the LHS $S_{T_n}^2 \to S_{T_{a,b}}^2$ a.s and $S_{T_n}^2 \leq \max(a^2, b^2)$ so by DCT we have $\lim_{n \to \infty} E[S_{T_n}^2] = E[S_{T_{a,b}}^2]$. Thus $E[T_{a,b}] = E[S_{T_{a,b}}^2]$.

    When $n \to \infty$, $\lim_{n \to \infty}$
    \begin{align*}
        E[S_{T_{n}}^2] = E[S_{T_{a,b}}^2] &= E[T_{a,b}] \\
        &= a^2 \cdot \frac{b}{a+b} + b^2 \cdot \frac{a}{a+b} \\
        &= ab
    \end{align*}
\end{example}
\begin{definition}[Walds Equation II for Random walks]
    If $\{X_n\}$ iid with mean $\mu$ $F_n = \sigma(X_1, \ldots, X_n)$ and $T$ is a stopping time wrt $\{F_n\}$ and $E[T] < \infty$

    Then $E[S_T] = \mu E[T]$. where $S_n = \sum_{i=1}^{n} X_i$.
    \begin{proof}
        Recall $S_n - n\mu$ is a martingale wrt $\{F_n\}$.
        \begin{align*}
            E[S_{n+1} - (n+1)\mu | F_n] &= E[S_n + X_{n+1} | F_n] - (n+1)\mu \\
            &= S_n + E[X_{n+1} | F_n] - (n+1)\mu \\
            &= S_n + \mu - (n+1)\mu \\
            &= S_n - n\mu
        \end{align*}
        If $E[S_T - T\mu] = 0$ then $E[S_T] = \mu E[T]$.
        Define $T_n = T \wedge n$ then by Walds I we have
        \begin{align*}
            E[S_{T_n} - T_n \mu] &= E[S_1 - \mu] = 0 \\
            E[S_{T_n}] &= \mu E[T_n]
        \end{align*}
        As $n \to \infty$, by MCT $E[T_n] \to E[T]$. For the LHS $S_{T_n} \to S_T$ a.s and $|S_{T_n}| \leq \sum_{i=1}^{T} |X_i|$.
        Since $E[T] < \infty$ and $E[|X_1|] < \infty$ we have $E[\sum_{i=1}^{T} |X_i|] = E[E[\sum_{i=1}^{T} |X_i| | T]] = E[T E[|X_1|]] = E[T] E[|X_1|] < \infty$. Thus by DCT we have $\lim_{n \to \infty} E[S_{T_n}] = E[S_T]$. Therefore $E[S_T] = \mu E[T]$.
    \end{proof}
\end{definition}
\begin{example}[HTTH coin flip]
    Flip a fair coin and stop until seeing the pattern $HTTH$. What is the expected stopping time? 

    Solution consider the following games:
    At every time $k$, a new gambler $G_k$ enters the game with 1 dollar. 
    They weill bet 1 on $X_s= H$ if correct, bet 2 on T, 4 on T, 8 on H. 
    $G_k$ will stop either losing a round or winning all 4 bets.

    Let $Y_n^{(k)}$ we the gross game of $G_k$ ar time $n$ then let $M_n = \sum_{k=1}^{n} Y_n^{(k) -1} = \sum_{k=1}^n Y_n^{(k)} - n$ be the total profic of all the gablers at time $n$. Then $M_n$ is a martingale with respect to the filtration $\mathcal{F}_n = \sigma(X_1, \ldots, X_n)$.
    Apply Wald on $M_n$ at time $T$ 

\end{example}



\end{document}